\documentclass[11pt,a4paper]{article}
\usepackage{amsmath,amssymb,amsfonts}
\usepackage{graphicx}
\usepackage[margin=1in]{geometry}
\usepackage{longtable}

\title{Supplementary Material:\\On the Polynomial $D^4 - 2D^3 + 2D + 1$}

\author{}

\date{\today}

\begin{document}

\maketitle

\tableofcontents
\newpage

%==============================================================================
\section{Complete Polynomial Evaluation}
%==============================================================================

\subsection{Definition}

The polynomial under consideration is:
\begin{equation}
P(D) = D^4 - 2D^3 + 2D + 1
\end{equation}

\subsection{Step-by-Step Computation}

Full computation for each integer $D \in [1, 10]$:

\begin{table}[h]
\centering
\begin{tabular}{|c|r|r|r|r|r|l|}
\hline
$D$ & $D^4$ & $-2D^3$ & $+2D$ & $+1$ & $P(D)$ & Prime? \\
\hline
1 & 1 & $-2$ & 2 & 1 & 2 & Yes \\
2 & 16 & $-16$ & 4 & 1 & 5 & Yes \\
3 & 81 & $-54$ & 6 & 1 & 34 & No: $2 \times 17$ \\
4 & 256 & $-128$ & 8 & 1 & 137 & Yes \\
5 & 625 & $-250$ & 10 & 1 & 386 & No: $2 \times 193$ \\
6 & 1296 & $-432$ & 12 & 1 & 877 & Yes \\
7 & 2401 & $-686$ & 14 & 1 & 1730 & No: $2 \times 5 \times 173$ \\
8 & 4096 & $-1024$ & 16 & 1 & 3089 & Yes \\
9 & 6561 & $-1458$ & 18 & 1 & 5122 & No: $2 \times 13 \times 197$ \\
10 & 10000 & $-2000$ & 20 & 1 & 8021 & No: $13 \times 617$ \\
\hline
\end{tabular}
\caption{Complete evaluation for $D \in [1,10]$. Each computation is verifiable by hand.}
\end{table}

\subsection{Verification of Key Result}

For $D = 4$:
\begin{align}
P(4) &= 4^4 - 2 \times 4^3 + 2 \times 4 + 1 \\
     &= 256 - 2 \times 64 + 8 + 1 \\
     &= 256 - 128 + 8 + 1 \\
     &= 137
\end{align}

This is exact arithmetic with no approximation.

%==============================================================================
\section{Primality Analysis}
%==============================================================================

\subsection{Primality Tests}

\begin{center}
\begin{tabular}{|c|c|l|}
\hline
$D$ & $P(D)$ & Primality Proof \\
\hline
2 & 5 & $5 = 5$ (prime, $\sqrt{5} < 3$) \\
4 & 137 & $137 / p \notin \mathbb{Z}$ for $p \in \{2,3,5,7,11\}$, $\sqrt{137} < 12$ \\
6 & 877 & $877 / p \notin \mathbb{Z}$ for $p \in \{2,3,5,7,11,13,17,19,23,29\}$, $\sqrt{877} < 30$ \\
8 & 3089 & $3089 / p \notin \mathbb{Z}$ for primes $p < 56$, $\sqrt{3089} < 56$ \\
\hline
\end{tabular}
\end{center}

\subsection{Composite Factorizations}

\begin{center}
\begin{tabular}{|c|c|l|}
\hline
$D$ & $P(D)$ & Factorization Verification \\
\hline
3 & 34 & $34 = 2 \times 17$ ($2 \times 17 = 34$ \checkmark) \\
5 & 386 & $386 = 2 \times 193$ ($2 \times 193 = 386$ \checkmark) \\
7 & 1730 & $1730 = 2 \times 865 = 2 \times 5 \times 173$ \checkmark \\
9 & 5122 & $5122 = 2 \times 2561 = 2 \times 13 \times 197$ \checkmark \\
10 & 8021 & $8021 = 13 \times 617$ ($13 \times 617 = 8021$ \checkmark) \\
\hline
\end{tabular}
\end{center}

\subsection{Pattern Summary}

\begin{itemize}
    \item Even $D \in \{2, 4, 6, 8\}$: All $P(D)$ are prime
    \item Even $D = 10$: $P(10) = 8021 = 13 \times 617$ (composite)
    \item Odd $D \geq 3$: All $P(D)$ are composite
\end{itemize}

The pattern ``$P(D)$ prime for even $D$'' fails at $D = 10$.

%==============================================================================
\section{Fibonacci Connection}
%==============================================================================

\subsection{Fibonacci Sequence}

The Fibonacci sequence $\{F_n\}$:
\[
1, 1, 2, 3, 5, 8, 13, 21, 34, 55, 89, 144, 233, \ldots
\]

\subsection{Matches}

\begin{center}
\begin{tabular}{|c|c|c|l|}
\hline
$D$ & $P(D)$ & Fibonacci? & If yes, which $F_n$? \\
\hline
1 & 2 & Yes & $F_3 = 2$ \\
2 & 5 & Yes & $F_5 = 5$ \\
3 & 34 & Yes & $F_9 = 34$ \\
4 & 137 & No & (between $F_{11}=89$ and $F_{12}=144$) \\
5 & 386 & No & \\
6 & 877 & No & \\
\hline
\end{tabular}
\end{center}

The Fibonacci connection holds for $D \in \{1, 2, 3\}$ only.

%==============================================================================
\section{Physical Constants}
%==============================================================================

\subsection{Fine Structure Constant}

\textbf{Definition}: $\alpha = e^2 / (4\pi\varepsilon_0 \hbar c)$

\textbf{CODATA 2022 value}:
\begin{equation}
\alpha^{-1} = 137.035\,999\,177(21)
\end{equation}

\textbf{Uncertainty}: $u_r = 1.5 \times 10^{-10}$ (relative)

\textbf{Source}: CODATA Task Group on Fundamental Constants, 2022 adjustment.

\subsection{Comparison with 137}

\begin{align}
\alpha^{-1} - 137 &= 0.035\,999\,177 \\
\frac{\alpha^{-1} - 137}{\alpha^{-1}} &= 2.63 \times 10^{-4} \\
\frac{\alpha^{-1} - 137}{137} &= 2.63 \times 10^{-4} = 0.0263\%
\end{align}

%==============================================================================
\section{QED Running Coupling}
%==============================================================================

\subsection{Theoretical Framework}

The QED beta function (one-loop):
\begin{equation}
\beta_{\text{QED}}(\alpha) = \frac{2\alpha^2}{3\pi} \sum_f N_c^f Q_f^2 \, \Theta(\mu - m_f)
\end{equation}

where $N_c^f$ is the color factor, $Q_f$ the electric charge, and $\Theta$ the Heaviside step function.

\subsection{Integrated Running}

\begin{equation}
\alpha^{-1}(\mu) = \alpha^{-1}(\mu_0) - \frac{2}{3\pi} \sum_f N_c^f Q_f^2 \ln\frac{\mu^2}{\max(m_f, \mu_0)^2}
\end{equation}

\subsection{Numerical Values}

\begin{center}
\begin{tabular}{|l|c|l|}
\hline
Scale $\mu$ & $\alpha^{-1}(\mu)$ & Source \\
\hline
$\mu \to 0$ (Thomson limit) & $137.036$ & Low-energy value \\
$\mu = m_e = 0.511$ MeV & $137.036$ & Running negligible below $m_e$ \\
$\mu = m_\mu = 105.7$ MeV & $\sim 135.9$ & Electron loop contribution \\
$\mu = m_\tau = 1.777$ GeV & $\sim 133.5$ & $e, \mu$ contributions \\
$\mu = M_Z = 91.2$ GeV & $127.95(1)$ & PDG 2024 \\
\hline
\end{tabular}
\end{center}

\textit{Note: Intermediate values are illustrative and scheme-dependent. The monotonicity of the running is the only ingredient required for the non-crossing conclusion.}

\subsection{Monotonicity}

$\alpha^{-1}(\mu)$ is monotonically decreasing for $\mu > m_e$:
\begin{equation}
\frac{d\alpha^{-1}}{d\ln\mu} = -\frac{2}{3\pi} \sum_f N_c^f Q_f^2 < 0
\end{equation}

\subsection{Behavior Below $m_e$}

For $\mu < m_e$:
\begin{itemize}
    \item No SM particles with mass $< m_e$ carry electric charge
    \item The electron contribution decouples: $\Theta(\mu - m_e) = 0$
    \item RG running is strongly suppressed
    \item Effective charge depends on scheme and low-energy definition, but monotonicity implies \textbf{no scale exists where $\alpha^{-1}(\mu) = 137$}
\end{itemize}

\subsection{Integer Crossing Analysis}

Question: Does $\alpha^{-1}(\mu) = 137$ for any $\mu$?

\begin{itemize}
    \item For $\mu \leq m_e$: $\alpha^{-1}(\mu) \approx 137.036 > 137$
    \item For $\mu > m_e$: $\alpha^{-1}(\mu) < 137.036$, decreasing
    \item Starting point: $\alpha^{-1}(m_e) = 137.036$
\end{itemize}

\textbf{Conclusion}: $\alpha^{-1}(\mu) > 137$ for $\mu \leq m_e$, and $\alpha^{-1}(\mu)$ decreases monotonically for $\mu > m_e$ starting from $137.036$. There is no scale where $\alpha^{-1}(\mu) = 137$ exactly within Standard Model QED.

%==============================================================================
\section{Exhaustive Polynomial Survey}
%==============================================================================

To assess uniqueness, we enumerate all monic degree-4 polynomials with small integer coefficients that evaluate to 137 at $D = 4$.

\subsection{Search Space Definition}

Polynomials of the form:
\begin{equation}
Q(D) = D^4 + aD^3 + bD^2 + cD + d
\end{equation}
with $a, b, c, d \in \mathbb{Z}$.

\subsection{Constraint at $D = 4$}

\begin{align}
Q(4) &= 256 + 64a + 16b + 4c + d = 137 \\
\therefore \quad 64a + 16b + 4c + d &= -119
\end{align}

\subsection{Sample Solutions (First 10 of 138)}

Representative solutions from the exhaustive enumeration:

\begin{center}
\begin{tabular}{|r|r|r|r|l|}
\hline
$a$ & $b$ & $c$ & $d$ & Polynomial \\
\hline
$-2$ & 0 & 2 & 1 & $D^4 - 2D^3 + 2D + 1$ \quad $\leftarrow$ \textit{this paper} \\
$-2$ & 0 & 3 & $-3$ & $D^4 - 2D^3 + 3D - 3$ \\
$-2$ & 0 & 4 & $-7$ & $D^4 - 2D^3 + 4D - 7$ \\
$-2$ & 0 & 1 & 5 & $D^4 - 2D^3 + D + 5$ \\
$-2$ & 0 & 0 & 9 & $D^4 - 2D^3 + 9$ \\
$-2$ & 1 & $-2$ & 1 & $D^4 - 2D^3 + D^2 - 2D + 1$ \\
$-2$ & 1 & $-1$ & $-3$ & $D^4 - 2D^3 + D^2 - D - 3$ \\
$-2$ & 1 & 0 & $-7$ & $D^4 - 2D^3 + D^2 - 7$ \\
$-2$ & $-1$ & 6 & 1 & $D^4 - 2D^3 - D^2 + 6D + 1$ \\
$-2$ & $-1$ & 7 & $-3$ & $D^4 - 2D^3 - D^2 + 7D - 3$ \\
\hline
\end{tabular}
\end{center}

\textit{The complete list of all 138 solutions is provided as supplementary data.}

\subsection{Count}

For each fixed $(a, b, c)$, there is exactly one $d$ satisfying the constraint:
\begin{equation}
d = -119 - 64a - 16b - 4c
\end{equation}

With $a, b, c \in [-10, 10]$ (21 choices each):
\begin{itemize}
    \item Total $(a, b, c)$ combinations: $21^3 = 9261$
    \item For each $(a,b,c)$, $d$ is uniquely determined by $d = -119 - 64a - 16b - 4c$
    \item The constraint $d \in [-10, 10]$ filters valid solutions
    \item \textbf{Exact count (by enumeration): 138 solutions}
\end{itemize}

\textbf{Conclusion}: Exhaustive non-uniqueness depends on the coefficient model:
\begin{itemize}
    \item \textbf{Monic-only} ($a_4 = 1$ fixed): \textbf{138} solutions
    \item \textbf{General quartic} ($a_4 \in [-10,10]$): \textbf{716} solutions
\end{itemize}

The polynomial $P(D) = D^4 - 2D^3 + 2D + 1$ is \textbf{not unique}; any physical interpretation must specify a principled mechanism selecting a particular coefficient model and polynomial.

\subsection{What Makes $P(D)$ Notable}

Among these polynomials, $P(D) = D^4 - 2D^3 + 2D + 1$ has:
\begin{itemize}
    \item Missing $D^2$ term (coefficient 0)
    \item Small coefficients: $\{-2, 0, 2, 1\}$
    \item Additional properties: primes at even $D \in \{2,4,6,8\}$, Fibonacci at $D \in \{2,3\}$
\end{itemize}

Whether these properties are significant or coincidental cannot be determined without a physical theory.

%==============================================================================
\section{Algebraic Properties of $P(D)$}
%==============================================================================

\subsection{Coefficient Structure}

\begin{equation}
P(D) = D^4 - 2D^3 + 0 \cdot D^2 + 2D + 1
\end{equation}

Coefficients: $(1, -2, 0, 2, 1)$

\subsection{Symmetry}

The coefficients are \textbf{not} palindromic: $(1, -2, 0, 2, 1) \neq (1, 2, 0, -2, 1)$.

\subsection{Irreducibility}

$P(D)$ is irreducible over $\mathbb{Q}$ (factorization over $\mathbb{Q}$ returns irreducible; verified by CAS).

\subsection{Roots}

Solving $P(D) = 0$ numerically (computed via NumPy):
\begin{align}
D_1 &= 1.5290855136 + 0.7429341359i \\
D_2 &= 1.5290855136 - 0.7429341359i \\
D_3 &= -0.5290855136 + 0.2570658641i \\
D_4 &= -0.5290855136 - 0.2570658641i
\end{align}

All four roots are complex (two conjugate pairs); no real roots exist.

%==============================================================================
\section{Historical Attempts at $\alpha^{-1}$}
%==============================================================================

\subsection{Eddington (1929--1944)}

Eddington proposed $\alpha^{-1} = 136$ based on counting degrees of freedom in a $16 \times 16$ matrix structure, later revised to 137 after experimental refinement.

\textbf{Status}: Rejected as numerology; no derivation from physical principles.

\subsection{Wyler (1969)}

Armand Wyler proposed:
\begin{equation}
\alpha = \frac{9}{8\pi^4} \left( \frac{\pi^5}{2^4 \cdot 5!} \right)^{1/4}
\end{equation}

Numerical value: $\alpha^{-1} \approx 137.03608$

\textbf{Status}: Remarkable numerical agreement, but the formula was not derived from any physical theory.

\subsection{Current Consensus}

The fine structure constant $\alpha$ is a free parameter of the Standard Model. Its value is determined by measurement, not calculation. Explanations include:
\begin{itemize}
    \item \textbf{Anthropic}: $\alpha$ takes a value compatible with complex chemistry and observers
    \item \textbf{String landscape}: $\alpha$ is fixed by vacuum selection in a multiverse
    \item \textbf{Unknown UV physics}: A deeper theory may predict $\alpha$
\end{itemize}

No first-principles derivation of $\alpha$ exists as of 2024.

\subsection{Methodological Distinction}

Unlike historical numerological fits, the present work:
\begin{itemize}
    \item Explicitly states non-uniqueness (many polynomials yield 137)
    \item Documents counterexamples ($D=10$ breaks the prime pattern)
    \item Provides falsifiable claims (RG crossing hypothesis → falsified)
    \item Separates mathematical identity from physical interpretation
\end{itemize}

The value lies not in claiming a ``derivation'' of $\alpha^{-1}$, but in demonstrating a systematic framework for evaluating such claims.

%% ============================================================================
%% APPENDIX: ALGEBRAIC EXHAUSTIVE EXPLORATION
%% ============================================================================

\section{Algebraic Exhaustive Exploration}

All computations in this section use Extended Precision Rational (EPR) arithmetic---exact rational numbers with no floating-point approximation.

\subsection{Exact Polynomial Evaluation}

The polynomial $P(D) = D^4 - 2D^3 + 2D + 1$ evaluated at $D = 4$:
\begin{align}
P(4) &= 4^4 - 2 \cdot 4^3 + 2 \cdot 4 + 1 \notag \\
     &= 256 - 128 + 8 + 1 \notag \\
     &= 137 \quad \text{(exact integer, algebraic identity)}
\end{align}

\subsection{Exact Difference from $\alpha^{-1}$}

Using the CODATA 2022 value $\alpha^{-1} = 137.035999177(21)$:
\begin{equation}
\alpha^{-1} - 137 = 0.035999177 \approx 0.036 \quad \text{(exact difference from CODATA 2022)}
\end{equation}

This confirms that 137 is \emph{not} equal to $\alpha^{-1}$; the coincidence is approximate.

\subsection{Complete $P(D)$ Evaluation Table}

\begin{center}
\begin{tabular}{cccc}
\hline
$D$ & $P(D)$ & Prime? & Factorization \\
\hline
1 & 2 & Yes & 2 \\
2 & 5 & Yes & 5 \\
3 & 34 & No & $2 \times 17$ \\
4 & 137 & Yes & 137 \\
5 & 386 & No & $2 \times 193$ \\
6 & 877 & Yes & 877 \\
7 & 1730 & No & $2 \times 5 \times 173$ \\
8 & 3089 & Yes & 3089 \\
9 & 5122 & No & $2 \times 13 \times 197$ \\
10 & 8021 & No & $13 \times 617$ \\
\hline
\end{tabular}
\end{center}

\textbf{Observation}: $P(D)$ is prime for $D \in \{1, 2, 4, 6, 8\}$ but composite for $D = 10$. The ``even $D$ yields prime'' hypothesis is \textbf{falsified} at $D = 10$.

\subsection{Exhaustive Polynomial Search}

For monic degree-4 polynomials $Q(D) = D^4 + aD^3 + bD^2 + cD + d$ with integer coefficients in $[-5, 5]$:
\begin{itemize}
    \item Total polynomials satisfying $Q(4) = 137$: \textbf{61}
    \item This confirms \textbf{non-uniqueness}: many polynomials produce the target value.
\end{itemize}

\subsection{Auto-Generated Physics Model Containers}

Based on the algebraic exploration, the following model containers were generated:

\begin{enumerate}
    \item \textbf{DimensionalCoupling}: Hypothesis $\alpha^{-1}(D) = D^4 - 2D^3 + 2D + 1$
    \begin{itemize}
        \item Evidence: $P(4) = 137 \approx \alpha^{-1}$ at $D = 4$
        \item Falsifiable: Measure $\alpha$ in effective theories with $D \neq 4$
    \end{itemize}
    
    \item \textbf{PrimeDimension}: Hypothesis ``stable physics requires $P(D)$ prime''
    \begin{itemize}
        \item Counter-evidence: $P(10) = 8021 = 13 \times 617$
        \item Status: \textbf{Falsified}
    \end{itemize}
    
    \item \textbf{KoideMassRelation}: Formula $Q = \frac{m_e + m_\mu + m_\tau}{(\sqrt{m_e} + \sqrt{m_\mu} + \sqrt{m_\tau})^2} = \frac{2}{3}$
    \begin{itemize}
        \item Status: Verified to 0.001\%
        \item Prediction: Extends to quarks? (testable)
    \end{itemize}
\end{enumerate}

\subsection{Summary of Verified and Falsified Claims}

\begin{center}
\begin{tabular}{lcc}
\hline
Claim & Status & Evidence \\
\hline
$P(4) = 137$ & \textbf{Verified} & Algebraic identity \\
Koide $Q = 2/3$ & \textbf{Verified} & Within 0.001\% \\
$P(D)$ prime for all even $D$ & \textbf{Falsified} & $P(10) = 13 \times 617$ \\
$P(D)$ unique for 137 & \textbf{Falsified} & 138 (monic) / 716 (general) alternatives \\
Multi-constant constraints & \textbf{Not found} & No formula constrains multiple constants \\
\hline
\end{tabular}
\end{center}

\end{document}
